% Gerado por tcc/analise/resultados_c3.py -- não edite à mão; rode o script de novo.
\begin{table}[htb]
\caption{Comparação entre os grupos por subperíodo}\label{tab:subperiodos}
\centering\footnotesize\setlength{\tabcolsep}{4pt}
\begin{tabular}{llrrcc}
\toprule
Indicador & Subperíodo & F & NF & p (t) & p (MW) \\
\midrule
Dív. bruta / Ativo & 2018--2019 & 25,4\% & 27,3\% & 0,815 & 0,625 \\
 & 2020--2021 & 27,8\% & 33,0\% & 0,543 & 0,255 \\
 & 2022--2025 & 25,5\% & 31,0\% & 0,497 & 0,416 \\
Exigível / Ativo & 2018--2019 & 53,1\% & 50,6\% & 0,703 & 0,704 \\
 & 2020--2021 & 56,8\% & 58,1\% & 0,825 & 0,871 \\
 & 2022--2025 & 55,8\% & 58,1\% & 0,760 & 0,957 \\
Dív. CP / Dív. total & 2018--2019 & 24,9\% & 33,4\% & 0,435 & 0,298 \\
 & 2020--2021 & 19,5\% & 36,4\% & \textbf{0,040}$^{**}$ & \textbf{0,043}$^{**}$ \\
 & 2022--2025 & 18,4\% & 41,5\% & 0,130 & 0,098$^{*}$ \\
\bottomrule
\end{tabular}
\fonte{Elaborado pelo autor (2026), com dados da Economática.}
\nota{Testes bicaudais sobre a média de cada empresa nas datas de 31/12 do período. p (MW): teste de Mann-Whitney, aproximação normal com correção de continuidade. $^{***}$, $^{**}$ e $^{*}$: significativo a 1\%, 5\% e 10\%. Subperíodos: pré-pandemia (2018--2019), pandemia (2020--2021) e pós-pandemia (2022--2025).}
\end{table}
